\section{Polyphase coordinates and folded bands}
\label{a-spectral-polyphase}

The finite Haar decomposition isolates an average and its new detail
channels. A standard polyphase decomposition instead groups the entries
of a cyclic lattice by their residue class. Both transformations preserve
the full state space. An elimination begins only when some coordinates
are solved for and removed. The spectral parameter then enters the
remaining operator, and the eliminated degrees of freedom return as
memory in time.

We use complex Hilbert spaces with inner product linear in the second
argument and set $\hbar=1$. All matrices in the finite sections act on
their stated coordinate spaces. In particular, the inverse of a low
channel block is a low channel operator, rather than a full space matrix.

\subsection{Residue coordinates and exact block reconstruction}

Fix an integer $p\geq2$ and $M\geq3$, and put $N=pM$. Primality is not
needed for the identities in this section. For $0\leq r<p$, define the
isometry $J_r:\C^M\to\C^N$ by
\[
 J_r|k\rangle=|pk+r\rangle,\qquad 0\leq k<M.
\]
Then $J_r^*J_s=\delta_{rs}I_M$ and $\sum_rJ_rJ_r^*=I_N$. The standard
analysis map $\mathcal A$ sends $|pk+r\rangle$ to
$|r\rangle\otimes|k\rangle$. This is a permutation unitary. For any
operator $X$ on $\C^N$, its residue blocks $X_{rs}=J_r^*XJ_s$ reconstruct it:
\begin{equation}
 X=\sum_{r,s=0}^{p-1}J_rX_{rs}J_s^*,\qquad
 (\mathcal A X\mathcal A^*)_{rs}=X_{rs}.
 \label{a-spectral-block-reconstruction}
\end{equation}
The first equality follows by inserting $\sum_rJ_rJ_r^*$ on both sides
of $X$. If $H=H^*$, its blocks satisfy $H_{rs}^*=H_{sr}$. If $U$ is
unitary, the same insertion gives
\begin{align}
 \sum_rU_{rs}^*U_{rt}&=J_s^*U^*UJ_t=\delta_{st}I_M,\nonumber\\
 \sum_sU_{rs}U_{ts}^*&=J_r^*UU^*J_t=\delta_{rt}I_M.
 \label{a-spectral-paraunitarity}
\end{align}
These are the block unitary identities. Every coordinate is retained.

For finite $H$, expansion of the matrix exponential proves
\[
 e^{-it\mathcal A H\mathcal A^*}
   =\mathcal A e^{-itH}\mathcal A^*.
\]
Indeed $(\mathcal A H\mathcal A^*)^n=\mathcal A H^n\mathcal A^*$,
and the series converges in matrix norm. The same transformation of
each Kraus operator gives a unitarily equivalent quantum channel:
if $\Phi(\rho)=\sum_\mu K_\mu\rho K_\mu^*$ with
$\sum_\mu K_\mu^*K_\mu=I$, then
\[
 K_\mu^{\rm poly}=\mathcal A K_\mu\mathcal A^*,\qquad
 \sum_\mu(K_\mu^{\rm poly})^*K_\mu^{\rm poly}=I.
\]
This is a change of coordinates in a given channel. It does not supply
a new dissipative dynamics.

\Needspace{16\baselineskip}
\begin{proposition}[Probability currents between residue channels]
\label{a-spectral-channel-current}
Suppose $i\dot\psi_r=\sum_sH_{rs}\psi_s$ with $H_{rs}^*=H_{sr}$.
The channel weights $w_r=\norm{\psi_r}^2$ satisfy
\[
 \dot w_r=\sum_s\mathcal J_{r\leftarrow s},\qquad
 \mathcal J_{r\leftarrow s}
   =2\operatorname{Im}\langle\psi_r,H_{rs}\psi_s\rangle,
 \qquad
 \mathcal J_{r\leftarrow s}=-\mathcal J_{s\leftarrow r}.
\]
Consequently $\sum_rw_r$ is constant.
\end{proposition}
\begin{proof}
Differentiating the norm gives
$\dot w_r=2\operatorname{Re}\langle\psi_r,-i\sum_sH_{rs}\psi_s\rangle$,
which is the stated sum of imaginary parts. Self-adjoint block symmetry
implies
$\langle\psi_s,H_{sr}\psi_r\rangle
=\overline{\langle\psi_r,H_{rs}\psi_s\rangle}$.
The currents therefore cancel in pairs, including a zero diagonal current.
\end{proof}
Here $w_r$ is a probability or squared amplitude weight. Calling it an
energy would require an additional energy observable and its balance law.

\subsection{A Bloch symbol with a fixed shift convention}

Let $S|n\rangle=|n+1\pmod N\rangle$ and
\begin{equation}
 H=\mu I+\kappa(S+S^*),\qquad \mu,\kappa\in\R.
 \label{a-spectral-free-cycle}
\end{equation}
Use the coarse Fourier matrix
$(F_M)_{jk}=M^{-1/2}e^{-2\pi ijk/M}$ and
$\theta_j=2\pi j/M$. After $\mathcal A$, Fourier transformation in $k$,
and the permutation from residue major to momentum major ordering, the
operator becomes a direct sum of $p\times p$ symbols. To see their entries,
note that $S$ sends $(r,k)$ to $(r+1,k)$ for $r<p-1$ and sends
$(p-1,k)$ to $(0,k+1)$. The coarse shift has Fourier eigenvalue
$e^{-i\theta_j}$. Thus define
\begin{equation}
 P_p(\theta)e_r=
 \begin{cases}
 e_{r+1},&r<p-1,\\
 e^{-i\theta}e_0,&r=p-1.
 \end{cases}
 \qquad
 \mathsf H_p(\theta)=\mu I_p+\kappa(P_p(\theta)+P_p(\theta)^*).
 \label{a-spectral-bloch-symbol}
\end{equation}
The transformed $H$ is $\bigoplus_{j=0}^{M-1}\mathsf H_p(\theta_j)$.
In particular, $P_p(\theta)^p=e^{-i\theta}I_p$.

\begin{proposition}[Exact folded cyclic spectrum]
\label{a-spectral-folded-bands}
For $q=0,\ldots,p-1$, the eigenvalues of
$\mathsf H_p(\theta)$ are
\begin{equation}
 E_q(\theta)=\mu+2\kappa
     \cos\!\left(\frac{\theta+2\pi q}{p}\right).
 \label{a-spectral-band-formula}
\end{equation}
At $\theta=\theta_j$, the union of these $pM$ eigenvalues, counting
multiplicity, is exactly the spectrum of \eqref{a-spectral-free-cycle}.
\end{proposition}
\begin{proof}
Put $\alpha_q=(\theta+2\pi q)/p$ and let the $r$th component of
$v_q$ be $p^{-1/2}e^{i\alpha_qr}$. On the interior components,
$P_p(\theta)v_q=e^{-i\alpha_q}v_q$. The wrapped component has the
same value because $e^{ip\alpha_q}=e^{i\theta}$. The $v_q$ are
orthonormal by the finite character sum. Applying
$\mu I+\kappa(P_p+P_p^*)$ gives \eqref{a-spectral-band-formula}.
Every fine Fourier index is uniquely $k=j+qM$. Its angle satisfies
$2\pi k/(pM)=(\theta_j+2\pi q)/p$. Substitution into the fine cyclic
eigenvalue $\mu+2\kappa\cos(2\pi k/N)$ proves the multiplicity statement.
\end{proof}

The same construction works for any circulant polynomial
$H=\sum_\ell c_\ell S^\ell$, with symbol
$\sum_\ell c_\ell P_p(\theta)^\ell$. It preserves all bands, whether or
not their energies coincide. For $N=p^h$, digit reversal relates this
standard residue split to the newest-digit adic split developed earlier.
It is still necessary to distinguish the retained residue-zero sites
from the retained Haar average channel. They are different subspaces.

\adicfigure{09_polyphase_symbol}{Three exact folded bands for the
nearest-neighbor cycle. The discrete markers represent the finite cyclic
spectrum. Band folding reorganizes its multiplicities without discarding
any eigenstate.}


\section{Regular Schur reduction and causal memory}
\label{a-spectral-feshbach}

Let a finite Hilbert space be the orthogonal sum
$\HH_L\oplus\HH_D$, and write a self-adjoint Hamiltonian as
\begin{equation}
 H=\begin{pmatrix}A&B\\B^*&D\end{pmatrix},\qquad A=A^*,\quad D=D^*.
 \label{a-spectral-block-hamiltonian}
\end{equation}
Here $B:\HH_D\to\HH_L$. Let $\iota_L$ embed $\HH_L$ into the full
space. This notation fixes the coordinate type of every compression.

\subsection{The resolvent chart and its exceptional energies}

\begin{theorem}[Finite Feshbach identity]
\label{a-spectral-schur-theorem}
If $z\notin\sigma(D)$, define
\begin{equation}
 \Sigma(z)=B(z-D)^{-1}B^*,\qquad
 \mathcal S(z)=z-A-\Sigma(z).
 \label{a-spectral-schur-definition}
\end{equation}
Then $z-H$ is invertible exactly when $\mathcal S(z)$ is invertible.
On this common regular set,
\begin{align}
 \iota_L^*(z-H)^{-1}\iota_L&=\mathcal S(z)^{-1},\nonumber\\
 \det(z-H)&=\det(z-D)\det\mathcal S(z).
 \label{a-spectral-schur-identities}
\end{align}
For $E\in\R\setminus\sigma(D)$, an eigenvector of $H$ is equivalent to
a nonzero solution of $\mathcal S(E)u=0$, with hidden component
$v=(E-D)^{-1}B^*u$.
\end{theorem}
\begin{proof}
Put $C=z-D$. Direct block multiplication gives
\[
 z-H=
 \begin{pmatrix}I&-BC^{-1}\\0&I\end{pmatrix}
 \begin{pmatrix}\mathcal S(z)&0\\0&C\end{pmatrix}
 \begin{pmatrix}I&0\\-C^{-1}B^*&I\end{pmatrix}.
\]
Both outer factors are invertible and have determinant one. This proves
the invertibility equivalence and determinant formula. Inverting the
factors gives the stated upper left block. Alternatively, in
$(z-H)(u,v)^{\mathsf T}=(f,0)^{\mathsf T}$, the lower equation gives
$v=C^{-1}B^*u$ and the upper equation becomes $\mathcal S(z)u=f$.
The same elimination with $f=0$ proves the eigenvector statement.
If $u=0$ at such an $E$, invertibility of $E-D$ forces $v=0$, so a
nonzero full eigenvector indeed has $u\ne0$.
\end{proof}

If $P_L=\iota_L\iota_L^*$, the full space compression is
$P_L(z-H)^{-1}P_L=\iota_L\mathcal S(z)^{-1}\iota_L^*$.
When a Haar transform $\mathcal W$ first exposes low and detail channels,
the embedding in the original coordinates is
$\iota_L=\mathcal W^*(|0\rangle\otimes I)$.
The inverse on $\HH_L$ must be embedded in this way before it is compared
with a full space matrix.

The determinant equality is initially a rational identity on the regular
set. Multiplication by $\det(z-D)$ cancels its denominators and yields the
full characteristic polynomial, which extends across $\sigma(D)$.
This extension does not make the displayed inverse of $z-D$ exist at
a pole. At a resonant energy the uneliminated equations determine which
hidden modes couple, survive, or cancel. A resonant subspace can be kept
as part of an enlarged visible block.

\subsection{Spectral residues, analytic sign, and branch weights}

Let $D=\sum_\lambda\lambda E_\lambda$ be its spectral resolution. Then
\begin{equation}
 \Sigma(z)=\sum_\lambda\frac{B E_\lambda B^*}{z-\lambda}.
 \label{a-spectral-self-energy-residues}
\end{equation}
Every residue $B E_\lambda B^*$ is positive semidefinite. Its rank records
which directions of the hidden eigenspace reach the visible space.
A vanishing residue removes the corresponding pole of $\Sigma$, even
though $z-D$ is still singular there.

\begin{proposition}[Upper half-plane sign and the exact visible norm]
\label{a-spectral-resolvent-sign}
For $z=x+iy$ with $y>0$,
\[
 \operatorname{Im}\Sigma(z)
  =-yB\bigl((x-D)^2+y^2I\bigr)^{-1}B^*\leq0,
 \qquad
 \operatorname{Im}\mathcal S(z)\geq yI.
\]
In particular $\mathcal S(z)$ is invertible and
$\norm{\mathcal S(z)^{-1}}\leq y^{-1}$.
For real $E\notin\sigma(D)$,
\[
 \Sigma'(E)=-B(E-D)^{-2}B^*.
\]
If $(u,v)$ is a normalized full eigenvector at $E$, then
\begin{equation}
 1=\langle u,(I-\Sigma'(E))u\rangle.
 \label{a-spectral-visible-norm}
\end{equation}
\end{proposition}
\begin{proof}
For real $\lambda$,
$\operatorname{Im}(z-\lambda)^{-1}=-y/((x-\lambda)^2+y^2)$.
Applying this identity to the spectral sum gives the first formula and
then the bound on $\operatorname{Im}\mathcal S$.
For any $u$, Cauchy--Schwarz implies
$y\norm{u}^2\leq|\langle u,\mathcal S(z)u\rangle|
\leq\norm{u}\norm{\mathcal S(z)u}$.
Thus $\mathcal S(z)$ has trivial kernel and an inverse of the stated
norm in finite dimension. Differentiating the resolvent gives the minus
sign in $\Sigma'$. The reconstructed hidden norm is
$\norm{v}^2=\langle u,B(E-D)^{-2}B^*u\rangle$.
Adding $\norm{u}^2$ proves \eqref{a-spectral-visible-norm}.
\end{proof}

For a one-dimensional visible block this says that an eigenstate's
visible probability is $[1-\Sigma'(E)]^{-1}$. It is at most one, as
required. Eliminating a normalized eigenvector therefore changes its
normalization metric as well as its spectral equation. Solving only a
nonlinear energy equation does not provide that probability weight.

\subsection{Exact elimination in the time domain}

\begin{theorem}[Causal memory with the initial hidden forcing]
\label{a-spectral-memory-theorem}
The finite system
$i\dot u=Au+Bv$, $i\dot v=Dv+B^*u$, with initial data $(u_0,v_0)$,
is equivalent to
\begin{align}
 i\dot u(t)&=Au(t)+B e^{-itD}v_0
              -i\int_0^tK(t-s)u(s)\,ds,\nonumber\\
 K(\tau)&=B e^{-i\tau D}B^*,\qquad \tau\geq0,
 \label{a-spectral-memory-equation}
\end{align}
together with reconstruction
\begin{equation}
 v(t)=e^{-itD}v_0-i\int_0^te^{-i(t-s)D}B^*u(s)\,ds.
 \label{a-spectral-hidden-reconstruction}
\end{equation}
For $\operatorname{Im}z>0$, the self-energy is its exact one-sided transform:
\begin{equation}
 \Sigma(z)=-i\int_0^\infty e^{iz\tau}K(\tau)\,d\tau.
 \label{a-spectral-memory-transform}
\end{equation}
\end{theorem}
\begin{proof}
Differentiate $e^{itD}v(t)$ and integrate the hidden equation. This gives
\eqref{a-spectral-hidden-reconstruction}, whose substitution in the
visible equation gives \eqref{a-spectral-memory-equation}. Conversely
the reconstruction differentiates to the hidden equation and supplies
the prescribed initial value, so it recovers the full evolution.
Since $D$ is self-adjoint, $\norm{K(\tau)}\leq\norm{B}^2$.
The factor $e^{-\operatorname{Im}z\,\tau}$ makes the transform absolutely
convergent. On a spectral eigenvalue of $D$,
$\int_0^\infty e^{i(z-\lambda)\tau}\,d\tau=i/(z-\lambda)$.
The finite spectral sum then proves \eqref{a-spectral-memory-transform}.
\end{proof}

The memory is causal because it uses only $s\leq t$. An initially empty
hidden sector removes the forcing $B e^{-itD}v_0$, not the convolution.
In finite dimension $K$ is a finite sum of oscillatory terms. Its decay,
a Markov approximation, or irreversible damping requires additional
limits or assumptions. The closed full dynamics preserves
$\norm{u(t)}^2+\norm{v(t)}^2$, while the visible norm alone can decrease
and later return. Energy-dependent Schur reduction is consequently
different from replacing $H$ by a fixed smaller self-adjoint matrix.


\section{Rooted trees and radial Jacobi response}
\label{a-spectral-radial}

Let $T_{p,h}$ be a rooted $p$-ary tree of depth $h\geq1$. Vertices at
depth $r$ have onsite energy $\epsilon_r\in\R$. Every parent-to-child
edge from level $r$ to $r+1$ has real coupling $\gamma_r/\sqrt p$,
where $\gamma_r\in\R$. The Hamiltonian is the corresponding real
symmetric adjacency matrix with onsite diagonal. The level symmetry
and this branch normalization are structural assumptions of the model.

\subsection{The radial reducing subspace}

Define normalized level states
\[
 |r\rangle=p^{-r/2}\sum_{|v|=r}|v\rangle,
 \qquad r=0,\ldots,h.
\]
They are orthonormal because different levels are disjoint and level
$r$ contains $p^r$ vertices. A direct action calculation gives
\begin{equation}
 H|r\rangle=\gamma_{r-1}|r-1\rangle
              +\epsilon_r|r\rangle+\gamma_r|r+1\rangle,
 \label{a-spectral-radial-action}
\end{equation}
with omitted terms at the endpoints. For example, the coefficient at
each child at level $r+1$ is
$p^{-r/2}\gamma_r/\sqrt p=\gamma_r p^{-(r+1)/2}$.
At each parent at level $r-1$, the $p$ child contributions instead give
$p\,p^{-r/2}\gamma_{r-1}/\sqrt p
=\gamma_{r-1}p^{-(r-1)/2}$.

Let $R:\C^{h+1}\to\ell^2(T_{p,h})$ send $e_r$ to $|r\rangle$.
Equation \eqref{a-spectral-radial-action} proves $HR=RJ_h$, where
\begin{equation}
 J_h=\begin{pmatrix}
 \epsilon_0&\gamma_0&&&\\
 \gamma_0&\epsilon_1&\gamma_1&&\\
 &\gamma_1&\epsilon_2&\ddots&\\
 &&\ddots&\ddots&\gamma_{h-1}\\
 &&&\gamma_{h-1}&\epsilon_h
 \end{pmatrix}.
 \label{a-spectral-jacobi-matrix}
\end{equation}
Because $H$ is self-adjoint, the orthogonal complement of $\Ran R$ is
also invariant: for $w\perp\Ran R$,
$\langle Re_r,Hw\rangle=\langle HRe_r,w\rangle=0$.
Hence this is a reducing decomposition, rather than only a compressed
approximation. The root belongs to the radial block.

\subsection{Angular chains and their multiplicities}

There is a more precise form of the complementary decomposition. For
each vertex $u$ at depth $a<h$, choose an orthonormal basis of vectors
$c=(c_0,\ldots,c_{p-1})\in\C^p$ with $\sum_dc_d=0$. For
$s=0,\ldots,h-a-1$, set
\begin{equation}
 |u,c;s\rangle
 =p^{-s/2}\sum_{d=0}^{p-1}c_d
       \sum_{w\in\{0,\ldots,p-1\}^{s}}|u d w\rangle.
 \label{a-spectral-angular-chain}
\end{equation}
For fixed $u,c$, these are unit vectors on successive levels. They obey
the Jacobi recurrence with onsite energies
$\epsilon_{a+1},\ldots,\epsilon_h$ and couplings
$\gamma_{a+1},\ldots,\gamma_{h-1}$. At $s=0$ the upward amplitude into
$u$ is proportional to $\sum_dc_d$, so it vanishes. At $s>0$ the same
branch-counting calculation as in \eqref{a-spectral-radial-action}
produces the neighboring chain vectors.

\begin{proposition}[Complete angular-chain decomposition]
\label{a-spectral-angular-completeness}
The radial chain and all the chains \eqref{a-spectral-angular-chain}
form an orthogonal decomposition of $\ell^2(T_{p,h})$. A chain starting
at depth $r$ has length $h-r+1$ and appears with multiplicity
$(p-1)p^{r-1}$, for $1\leq r\leq h$.
\end{proposition}
\begin{proof}
Within one chain, different $s$ have disjoint level supports. Chains
based at the same $u$ and different contrast vectors have inner product
$\langle c,c'\rangle=0$ at a common level. If two vertices are in disjoint
subtrees, their chains have disjoint support. If one vertex lies below
another, the higher chain is constant, up to its fixed ancestral
coefficient, on all children below the lower vertex at a common level.
The lower chain has zero sum over its own child contrasts, so their
inner product is zero. The same zero-sum calculation makes every
angular vector orthogonal to the uniform radial level vector.
At a fixed depth $b$, the radial vector contributes one dimension.
The angular vectors present there contribute
$\sum_{r=1}^b(p-1)p^{r-1}=p^b-1$ dimensions. They are orthogonal and
therefore form a basis of that $p^b$-dimensional level space. Taking
all levels proves completeness. There are $p^{r-1}$ choices of the
vertex $u$ at depth $r-1$ and $p-1$ contrasts at each, giving the
stated multiplicity.
\end{proof}

The rank $(p-1)p^{r-1}$ thus counts chains born at level $r$ or
new Haar directions on a level. A single level's detail space is
usually transported to the next level; it is the entire angular
chain that is invariant under the tree Hamiltonian. The full dark
dimension is
\[
 \sum_{r=1}^h(p-1)p^{r-1}(h-r+1)
 =\sum_{b=0}^hp^b-(h+1).
\]
Every one of these states is invisible to a root-localized symmetric
source because it is orthogonal to the radial reducing block.

\subsection{Continued fractions and their meromorphic reading}

For a tail $J_{r:h}$, put
$g_r(z)=\langle r,(z-J_{r:h})^{-1}r\rangle$. On a common regular
domain, scalar Schur elimination gives
\begin{align}
 g_h(z)&=(z-\epsilon_h)^{-1},\nonumber\\
 g_r(z)&=\frac1{z-\epsilon_r-\gamma_r^2g_{r+1}(z)},
          \qquad r<h.
 \label{a-spectral-continued-fraction}
\end{align}
Defining $\Sigma_h=0$ and
$\Sigma_r=\gamma_r^2/(z-\epsilon_{r+1}-\Sigma_{r+1})$ gives
\begin{equation}
 \langle\mathrm{root},(z-H)^{-1}\mathrm{root}\rangle
   =g_0(z)=\frac1{z-\epsilon_0-\Sigma_0(z)}.
 \label{a-spectral-root-green}
\end{equation}
The equality to the tree resolvent follows from the exact reducing
decomposition. Equations \eqref{a-spectral-continued-fraction} are
rational identities. Nested denominators can become singular at a
real energy where a canceled final rational expression is finite.
One should use the resolvent or the canceled rational function at
such a point, rather than perform division by zero in the nested chart.

The Jacobi coefficients in \eqref{a-spectral-jacobi-matrix} do not
depend on $p$. Thus the root Green function and radial propagator
are prime independent when the same $\epsilon_r,\gamma_r$ are used.
The complete tree spectrum retains the $p$-dependent angular
multiplicities. If the per-edge coupling is instead a fixed $g_r$,
the radial coupling is $\sqrt p\,g_r$ and its response depends on $p$.
Prime independence is therefore a consequence of the declared
normalization, not a universality statement about arbitrary trees.

\subsection{A visible two-component block}

As a conditional finite model, take
$A=H_D(k)=k\sigma_z+m_0\sigma_x$ with $k,m_0\in\R$ and
$D=J_{1:h}\otimes I_2$. Couple the visible spinor identically to the
first site of the tail by $B=\gamma_0(\langle1|\otimes I_2)$.
Its self-energy is scalar in spin space:
\[
 \Sigma_D(z)=\sigma(z)I_2,\qquad
 \sigma(z)=\gamma_0^2\langle1,(z-J_{1:h})^{-1}1\rangle.
\]
The determinant formula from Theorem~\ref{a-spectral-schur-theorem} is
\begin{equation}
 \det(z-H_{\rm full})
 =\det(z-D)\bigl((z-\sigma(z))^2-k^2-m_0^2\bigr).
 \label{a-spectral-visible-determinant}
\end{equation}
Consequently any eigenvalue in the regular chart obeys
\begin{equation}
 (E-\sigma(E))^2=k^2+m_0^2.
 \label{a-spectral-visible-branch}
\end{equation}
The formula is a dressed dispersion relation of this specific model.
Its mass, couplings, and hidden sector have been supplied. No physical
interaction or measured dispersion follows solely from the tree split.

For a nondegenerate visible spin direction with normalized low-space
direction $u_*$, the visible weight of its normalized full eigenstate is
$[1-\sigma'(E)]^{-1}$. If $k=m_0=0$, the spin eigenspace is degenerate,
but the same normalization relation holds for each chosen unit vector
in that eigenspace.

If every internal tail hopping is nonzero, the first tail site is cyclic.
One can see this without an abstract cyclicity theorem: a tail eigenvector
with zero first component has zero second component by the first row,
and successive rows force all later components to vanish. Therefore
every tail eigenvector has nonzero first component. At an eigenvalue
$\lambda$ of that tail, solvability of
$(\lambda-D)v=B^*u$ requires its right side to be orthogonal to the tail
eigenspace. With $\gamma_0\ne0$ and identical spin coupling this forces
$u=0$. The visible equation would then require $Bv=0$, which forces
$v=0$ on the tail eigenspace. Thus in this irreducible case the full
coupled spectrum avoids the tail eigenvalues. With a vanishing hopping,
coincident visible and hidden energies can occur and the regular-chart
qualification in \eqref{a-spectral-visible-branch} is essential.

The scalar causal kernel is
$K(\tau)=\gamma_0^2\langle1,e^{-i\tau J_{1:h}}1\rangle$, multiplying
$I_2$. It uses the same radial coefficients as the rational self-energy.
Its exact causal relation is \eqref{a-spectral-memory-transform}.
\adicfigure{08_radial_self_energy}{A branch-normalized finite tree has
the same root response as its radial Jacobi chain. The complex-energy
resolvent avoids real poles. The oscillatory memory is the return
amplitude of a finite hidden chain.}


\section{Chebyshev decimation and resonant channels}
\label{a-spectral-chebyshev}

We now eliminate intermediate residue sites in
\eqref{a-spectral-free-cycle}, assuming $\kappa\ne0$. The retained
sites are $C=\{0,p,\ldots,(M-1)p\}$ and the eliminated sites comprise
the $p-1$ interior sites between adjacent retained sites. Identify
$\ell^2(C)$ with $\C^M$ and let $T|j\rangle=|j+1\pmod M\rangle$.
Set
\begin{equation}
 x=\frac{z-\mu}{\kappa},\qquad
 R_n(x)=2T_n(x/2),
 \label{a-spectral-normalized-energy}
\end{equation}
where $T_n$ is the first-kind Chebyshev polynomial. The second-kind
polynomials $U_n$ satisfy $U_{-1}=0$, $U_0=1$ and
$U_{n+1}(t)=2tU_n(t)-U_{n-1}(t)$.

\subsection{The eliminated path fixes the Schur prefactor}

The stationary equation on an interior site is
$\psi_{n+1}=x\psi_n-\psi_{n-1}$. The eliminated block of $z-H$ in
one cell is $\kappa K_{p-1}(x)$, where $K_n(x)$ has diagonal $x$
and neighboring off-diagonal entries $-1$.

\begin{lemma}[Dirichlet path determinant and endpoint entries]
\label{a-spectral-path-inverse}
For $n\geq1$, $\det K_n(x)=U_n(x/2)$. If this determinant is nonzero,
\[
 (K_n(x)^{-1})_{11}=(K_n(x)^{-1})_{nn}
       =\frac{U_{n-1}(x/2)}{U_n(x/2)},\qquad
 (K_n(x)^{-1})_{1n}=\frac1{U_n(x/2)}.
\]
For $n=1$ these entries all denote the same value $1/x$.
\end{lemma}
\begin{proof}
Expansion along the last row gives the determinant recurrence
$d_n=xd_{n-1}-d_{n-2}$, with $d_0=1$ and $d_1=x$.
This is exactly the recurrence and initial data for $U_n(x/2)$.
Deleting the first or last row and column leaves $K_{n-1}$, proving
the diagonal inverse entries by cofactors. For the endpoint entry,
the cofactor consists of the product of $n-1$ entries $-1$ along the
unique path from the first endpoint to the last. Its cofactor sign is
$(-1)^{n+1}$, so the product has total sign
$(-1)^{n+1}(-1)^{n-1}=1$. Division by $d_n$ gives the stated value.
\end{proof}

\begin{theorem}[Exact free-cycle decimation]
\label{a-spectral-chebyshev-schur}
For $U_{p-1}(x/2)\ne0$, the Schur operator on the retained sublattice is
\begin{equation}
 \mathcal S_p(z)=\frac{\kappa}{U_{p-1}(x/2)}
          \bigl[R_p(x)I_M-(T+T^*)\bigr].
 \label{a-spectral-decimation-formula}
\end{equation}
Equivalently, $\mathcal S_p(z)=zI-H_{\rm eff}^{(p)}(z)$, where
\begin{align}
 H_{\rm eff}^{(p)}(z)&=\mu_p(z)I+\kappa_p(z)(T+T^*),\nonumber\\
 \kappa_p(z)&=\frac{\kappa}{U_{p-1}(x/2)},\qquad
 \mu_p(z)=z-\frac{\kappa R_p(x)}{U_{p-1}(x/2)}.
 \label{a-spectral-effective-coefficients}
\end{align}
\end{theorem}
\begin{proof}
The eliminated sites split into $M$ disjoint Dirichlet paths of length
$p-1$. Hence their resolvent exists exactly when
$U_{p-1}(x/2)\ne0$. A retained site touches one endpoint on each of
its two neighboring paths. Each path contributes the onsite
self-energy $\kappa U_{p-2}/U_{p-1}$ by
Lemma~\ref{a-spectral-path-inverse}; the two contributions add.
The path joining two adjacent retained sites contributes off-diagonal
self-energy $\kappa/U_{p-1}$. Thus
\[
 \mathcal S_p(z)
 =\kappa\left(x-\frac{2U_{p-2}(x/2)}{U_{p-1}(x/2)}\right)I
       -\frac{\kappa}{U_{p-1}(x/2)}(T+T^*).
\]
The identity
$xU_{p-1}(x/2)-2U_{p-2}(x/2)=2T_p(x/2)$ follows from the
three-term recurrence and $U_p-U_{p-2}=2T_p$.
For the latter identity, substituting $x/2=\cos\theta$ gives
$[\sin((p+1)\theta)-\sin((p-1)\theta)]/\sin\theta
=2\cos(p\theta)$, and polynomial equality extends it to all $x$.
This proves \eqref{a-spectral-decimation-formula} including its
prefactor. Rearranging its diagonal gives
\eqref{a-spectral-effective-coefficients}.
\end{proof}

The normalized effective energy is
$(z-\mu_p(z))/\kappa_p(z)=R_p(x)$. The parameter $z$ remains fixed
through the elimination. Since $\mu_p$ and $\kappa_p$ depend on it,
this is an energy-dependent spectral closure, not an energy-independent
Hamiltonian flow. For $p=2$, $R_2(x)=x^2-2$ and $U_1(x/2)=x$, giving
\begin{equation}
 H_{\rm eff}^{(2)}(z)
 =\mu I+\frac{\kappa^2}{z-\mu}(2I+T+T^*),\qquad z\ne\mu.
 \label{a-spectral-dyadic-effective}
\end{equation}
This decimation retains every other site. It does not retain the
average of each neighboring pair.

\subsection{The poles have an exact Dirichlet interpretation}

The sine formula $U_{p-1}(\cos\theta)=\sin(p\theta)/\sin\theta$
shows that the excluded normalized energies are
\begin{equation}
 x_j=2\cos(j\pi/p),\qquad
 z_j=\mu+2\kappa\cos(j\pi/p),\quad 1\leq j<p.
 \label{a-spectral-dirichlet-resonances}
\end{equation}
These are the energies of one eliminated path. Indeed its $j$th
eigenvector has interior components $\sin(j\pi r/p)$,
$r=1,\ldots,p-1$, and the nearest-neighbor recurrence gives precisely
$z_j$. Both boundary values at $r=0,p$ vanish. The eigenvectors are
orthogonal and there are $p-1$ of them, so these exhaust the path spectrum.

The polynomial map $R_p$ is defined at every $x_j$, while the Schur
coefficients in \eqref{a-spectral-effective-coefficients} are not.
Cancelling denominators in a determinant can cross this chart boundary;
dividing by them in an eigenvector reconstruction cannot. A correct
description there keeps the resonant cell modes or solves the full
finite system. The fixed poles of the eliminated path are not all
automatically eigenvalues of the full cyclic Hamiltonian.

Let $Q$ project onto the eliminated sites. On the regular chart,
\[
 \det(z-H)=\det(z-QHQ)\det\mathcal S_p(z),\qquad
 \det(z-QHQ)=\bigl[\kappa^{p-1}U_{p-1}(x/2)\bigr]^M.
\]
Their product is a polynomial identity after cancellation. It records
both the hidden path factor and the coarse Schur factor, including
their cancellations at the resonant energies.

\subsection{Transfer coordinates and many-to-one branch maps}

Define $A(x)=\begin{pmatrix}x&-1\\1&0\end{pmatrix}$. Its determinant
is one and induction gives
\begin{equation}
 A(x)^n=\begin{pmatrix}
 U_n(x/2)&-U_{n-1}(x/2)\\
 U_{n-1}(x/2)&-U_{n-2}(x/2)
 \end{pmatrix},\quad n\geq1.
 \label{a-spectral-transfer-power}
\end{equation}
For $n=1$ this uses $U_{-1}=0$. Multiplication by $A(x)$ advances
the first column through the Chebyshev recurrence and places the
negative previous first column in the second. The trace is $R_n(x)$.
Sampling the fine recurrence every $p$ steps therefore uses the
unimodular transfer matrix $A(x)^p$. Cayley--Hamilton gives the coarse
recurrence $\chi_{j+1}-R_p(x)\chi_j+\chi_{j-1}=0$ for
$\chi_j=\psi_{pj}$. It explains the normalized spectral polynomial.
The path-inverse proof above additionally supplies the prefactor needed
for the resolvent.

The folded branches from Proposition~\ref{a-spectral-folded-bands},
in normalized coordinates, satisfy
\begin{equation}
 x_{p,q}(\theta)=2\cos\!\left(\frac{\theta+2\pi q}{p}\right),
 \qquad R_p(x_{p,q}(\theta))=2\cos\theta.
 \label{a-spectral-branch-relation}
\end{equation}
The second equality is just $T_p(\cos\alpha)=\cos(p\alpha)$.
For a coarse normalized energy in the open interval $(-2,2)$,
choose its angle in $(0,\pi)$. These formulas supply the $p$ distinct
real roots of $R_p(x)=2\cos\theta$. At the endpoints, some roots
coincide and algebraic multiplicities must be retained. Globally the
branches exchange labels when the angle crosses a cell boundary.
There is no single-valued inverse of $R_p$ on $[-2,2]$.

For the fine cycle, $x_k=2\cos(2\pi k/(pM))$ maps to
$2\cos(2\pi k/M)$. Each coarse Fourier label has exactly $p$ fine
preimages. If two labels already have the same numerical energy,
their multiplicities add as usual. Band folding retains the branch
label and its vector, whereas normalized decimation identifies all
of these labels at the same coarse energy. Their spectral relation
does not make elimination an invertible operator change of basis.


\section{Composition, transfer traces, and the free spectral measure}
\label{a-spectral-semigroup}

The spectral maps compose more simply than the energy-dependent
operators that produce them. Keeping the normalized polynomial map,
the rational Schur chart, and the transfer matrix distinct prevents an
exact free-chain identity from being used outside its model.

\subsection{Polynomial composition and iterated regular elimination}

\begin{theorem}[Chebyshev composition and Schur composition]
\label{a-spectral-semigroup-theorem}
For all positive integers $m,n$, the polynomial identities
\begin{align}
 R_n(R_m(x))&=R_{mn}(x)=R_m(R_n(x)),\nonumber\\
 U_{mn-1}(t)&=U_{n-1}(T_m(t))\,U_{m-1}(t)
 \label{a-spectral-polynomial-semigroup}
\end{align}
hold over $\C$. In a free cycle with $N=mnL$, $L\geq3$, successive
$m$- then $n$-site decimation at a fixed $z$ agrees with direct
$mn$-site decimation whenever its intermediate Schur inverses exist.
Its final coefficients are
\begin{equation}
 \kappa_{mn}(z)=\frac{\kappa}{U_{mn-1}(x/2)},\qquad
 \mu_{mn}(z)=z-\frac{\kappa R_{mn}(x)}{U_{mn-1}(x/2)}.
 \label{a-spectral-iterated-coefficients}
\end{equation}
\end{theorem}
\begin{proof}
For $x=2\cos\theta$, first-kind composition gives
$R_n(R_m(x))=2\cos(mn\theta)=R_{mn}(x)$.
It holds on an interval and is therefore an identity of polynomials,
including complex $x$. For $t=\cos\theta$ away from zeros of the sine
denominators,
\[
 U_{n-1}(T_m(t))U_{m-1}(t)
 =\frac{\sin(mn\theta)}{\sin(m\theta)}
       \frac{\sin(m\theta)}{\sin\theta}
 =U_{mn-1}(t).
\]
Polynomial equality extends it through those zeros and to all $t$.
After the first regular elimination, the normalized variable is
$x_m=(z-\mu_m)/\kappa_m=R_m(x)$. A second one gives
\[
 \kappa_{m,n}
 =\frac{\kappa_m}{U_{n-1}(x_m/2)}
 =\frac{\kappa}
   {U_{m-1}(x/2)U_{n-1}(T_m(x/2))}
 =\frac{\kappa}{U_{mn-1}(x/2)}.
\]
Similarly
$\mu_{m,n}=z-\kappa_{m,n}R_n(R_m(x))$, which gives
\eqref{a-spectral-iterated-coefficients}.
At each stage the matrices remain scalar nearest-neighbor circulants
with the displayed complex coefficients, so their algebraic Schur
formula continues to apply. These are rational identities on a common
regular chart. They need not be interpreted as self-adjoint matrices at
complex $z$.
\end{proof}

The product identity in \eqref{a-spectral-polynomial-semigroup} makes
the exceptional denominators explicit. A nonzero final denominator
implies both denominators along that two-stage chart are nonzero.
At a zero of the final denominator the polynomial trace composition
is still valid, but the divided Schur formulas require a new chart or
retained resonant modes. Iterating one prime $p$ through depth $h$,
where the cycle has the required divisibility, replaces $p$ by $p^h$
in \eqref{a-spectral-iterated-coefficients}. Different prime orders
give the same regular free-chain result. A nontranslation-invariant
onsite potential, an interacting Hamiltonian, or a field-dependent
constraint elimination needs its own closure analysis.

\subsection{What remains true for a general unimodular transfer matrix}

Let $M_0\in SL_2(\C)$ and $x=\operatorname{tr}M_0$. No
diagonalizability assumption is needed. Cayley--Hamilton gives
$M_0^2-xM_0+I=0$, hence
\begin{equation}
 M_0^n=U_{n-1}(x/2)M_0-U_{n-2}(x/2)I,\qquad n\geq1,
 \label{a-spectral-sl2-power}
\end{equation}
with $U_{-1}=0$. For $n=1$ it is the identity $M_0=M_0$, and for
$n=2$ it is Cayley--Hamilton. Multiplying by $M_0$ and replacing
$M_0^2$ by $xM_0-I$ proves the induction step. Taking traces yields
\begin{equation}
 \operatorname{tr}(M_0^n)=R_n(\operatorname{tr}M_0).
 \label{a-spectral-sl2-trace}
\end{equation}
The coefficient identity is the same
$xU_{n-1}-2U_{n-2}=2T_n$ used above.

If $w_{k+1}=M_0w_k$, each residue subsequence satisfies
\[
 w_{nq+j}=M_0^j(M_0^n)^q w_0,\qquad 0\leq j<n.
\]
Thus ordinary matrix powering produces the normalized trace semigroup,
including matrices with repeated eigenvalues at trace $\pm2$.
The trace alone does not determine $M_0$, its eigenvectors, or a
local Hamiltonian realization. In the free chain, $M_0=A(x)$ has
the required nearest-neighbor realization. For a general transfer
problem only the matrix power and trace statements have been proved.

\subsection{The arcsine measure and the Koopman operator}

On $[-2,2]$ define the probability measure
\begin{equation}
 d\nu(x)=\frac{dx}{\pi\sqrt{4-x^2}}.
 \label{a-spectral-arcsine-measure}
\end{equation}
It is the distribution of $2\cos\theta$ for a uniform angle on the
circle, or on $[0,\pi]$. It is also the limiting normalized spectral
measure of the free cycles: for a continuous test function $f$,
\[
 \frac1N\sum_{k=0}^{N-1}f(2\cos(2\pi k/N))
 \longrightarrow\frac1{2\pi}\int_0^{2\pi}f(2\cos\theta)\,d\theta.
\]
This is the Riemann-sum limit of a continuous periodic function.

\Needspace{17\baselineskip}
\begin{proposition}[Invariant free measure and exact mode routing]
\label{a-spectral-arcsine-invariance}
For every integer $p\geq2$, $(R_p)_*\nu=\nu$. The map
$\mathcal K_pf=f\circ R_p$ is an isometry on $L^2(\nu)$. With
\[
 \phi_0=1,\qquad \phi_n(x)=\sqrt2\,T_n(x/2),\quad n\geq1,
\]
its action and adjoint on an orthonormal basis are
\begin{equation}
 \mathcal K_p\phi_n=\phi_{pn},\qquad
 \mathcal K_p^*\phi_m=
 \begin{cases}\phi_{m/p},&p\mid m,\\0,&p\nmid m.\end{cases}
 \label{a-spectral-koopman-routing}
\end{equation}
In the first case $m=0$ is included, so $\mathcal K_p^*\phi_0=\phi_0$.
\end{proposition}
\begin{proof}
For a bounded measurable $f$, use the full circle parameterization:
\begin{align*}
 \int f(R_p(x))\,d\nu(x)
 &=\frac1{2\pi}\int_0^{2\pi}f(2\cos(p\theta))\,d\theta\\
 &=\frac1{2\pi p}\int_0^{2\pi p}f(2\cos u)\,du
 =\int f(x)\,d\nu(x).
\end{align*}
The last equality uses periodicity. Applying this equality to $|f|^2$
proves the isometry, first on bounded functions and then by truncation
on all of $L^2(\nu)$. Pulling back to $[0,\pi]$ makes the $\phi_n$
the normalized cosine functions, and elementary cosine orthogonality
proves their orthonormality. Their span is dense: Chebyshev polynomials
span ordinary polynomials; those approximate continuous functions
uniformly, for example by Bernstein polynomials after rescaling to
$[0,1]$. The approximation follows from uniform continuity and the
binomial variance bound, by separating $|k/n-t|\leq\delta$ from its
complement and then letting $n\to\infty$, $\delta\to0$.
Continuous functions are dense for this finite measure, as one sees by
approximating interval steps in the angle coordinate by continuous
ramps. This supplies completeness.
Finally, Chebyshev composition gives $\mathcal K_p\phi_n=\phi_{pn}$.
For every $n,m\geq0$,
$\langle\phi_n,\mathcal K_p^*\phi_m\rangle
=\langle\phi_{pn},\phi_m\rangle=\delta_{pn,m}$.
Completeness yields the adjoint formula.
\end{proof}

$\mathcal K_p$ is an isometry with proper range, rather than a unitary:
for example $\phi_1$ is orthogonal to its range. Its adjoint discards
nondivisible modes. Invariance of $\nu$ is a statement about the
normalized free spectral distribution under $R_p$. It is not a proof
that a sequence of growing Hilbert spaces has converged to an interacting
Hamiltonian or that its quantum state converges. The direct-limit
operator construction has separate hypotheses in the final part.

The mean logarithmic expansion of this map is also exact:
\begin{equation}
 \int_{-2}^2\log|R_p'(x)|\,d\nu(x)=\log p.
 \label{a-spectral-log-expansion}
\end{equation}
Indeed for $x=2\cos\theta$,
$R_p'(x)=p\sin(p\theta)/\sin\theta$. Logarithms at its finitely
many zeros are integrable, since they have the local form
$\log|\theta-\theta_0|$ plus a bounded term. Hence the integral is
$\log p$ plus the difference of the integrals of
$\log|\sin(p\theta)|$ and $\log\sin\theta$ over $[0,\pi]$.
The substitution $u=p\theta$ and $\pi$-periodicity of
$\log|\sin u|$ make these two integrals equal. This calculation
does not assert the entropy of a different physical dynamics.

\adicfigure{11_prime_renormalization}{The normalized maps
$R_2$, $R_3$, and $R_5$ preserve the same arcsine spectral measure.
The accompanying bounded moment checks illustrate the exact measure
identity, while each Schur realization retains its own excluded poles.}


\section{Analytic residue lanes and hypergeometric functions}
\label{a-spectral-hypergeometric}

Residue filtering also acts on coefficients of a convergent analytic
series. Here the input is a function, not a quantum state, and the
projector acts on its coefficient sequence. Evaluating one resulting
lane at an operator does not generally turn that lane into a Hilbert
space projector.

\subsection{Character filtering and factorial separation}

Fix an integer $L\geq2$, write $\omega_L=e^{2\pi i/L}$, and let
$F(z)=\sum_{n=0}^\infty a_nz^n$ be analytic for $|z|<\rho$.
The whole disk is invariant under the rotations used below. Define
\begin{equation}
 (\mathfrak R_{L,j}F)(z)
 =\frac1L\sum_{\ell=0}^{L-1}\omega_L^{-j\ell}F(\omega_L^\ell z),
 \quad 0\leq j<L.
 \label{a-spectral-analytic-filter}
\end{equation}
Absolute convergence for $|z|<\rho$ and the finite character sum give
\begin{equation}
 (\mathfrak R_{L,j}F)(z)=\sum_{q=0}^\infty a_{Lq+j}z^{Lq+j},
 \quad
 \mathfrak R_{L,j}\mathfrak R_{L,k}
   =\delta_{jk}\mathfrak R_{L,j},\quad
 \sum_j\mathfrak R_{L,j}=I.
 \label{a-spectral-analytic-lanes}
\end{equation}
To prove the coefficient formula, the coefficient of $a_nz^n$ in
\eqref{a-spectral-analytic-filter} is
$L^{-1}\sum_\ell\omega_L^{(n-j)\ell}$, equal to one exactly when
$n\equiv j\pmod L$. The idempotent and reconstruction laws follow
by filtering these coefficients again. For an entire $F$ the formulas
hold everywhere. For a function continued outside its disk, the
domain and any chosen branches must also support all rotations; the
local power-series identity does not remove such analytic restrictions.

Apply this filter to $e^z$ and put
\begin{equation}
 E_{L,j}(z)=\sum_{q=0}^\infty\frac{z^{Lq+j}}{(Lq+j)!}.
 \label{a-spectral-factorial-lane}
\end{equation}
For $\alpha\in\C$, write
$(\alpha)_q=\alpha(\alpha+1)\cdots(\alpha+q-1)$ and
$(\alpha)_0=1$. Splitting the factorial into step-$L$ factors gives
\begin{align}
 \frac{(Lq+j)!}{j!}
 &=\prod_{r=1}^{L}\prod_{k=0}^{q-1}(j+r+Lk)\nonumber\\
 &=L^{Lq}\prod_{r=1}^{L}\left(\frac{j+r}{L}\right)_q.
 \label{a-spectral-factorial-split}
\end{align}
Exactly one bottom parameter is one, at $r=L-j$. Define the multiset
\[
 \mathcal A_{L,j}
 =\left\{\frac{j+r}{L}:1\leq r\leq L\right\}\setminus\{1\}.
\]
Every member is positive, so none is a forbidden nonpositive integer
in a hypergeometric denominator. Since $(1)_q=q!$,
\begin{equation}
 E_{L,j}(z)=\frac{z^j}{j!}\,
 {}_0F_{L-1}\!\left(\begin{matrix}-\\\mathcal A_{L,j}\end{matrix};
                              \frac{z^L}{L^L}\right).
 \label{a-spectral-hypergeometric-lane}
\end{equation}
Here ${}_0F_s(;b_1,\ldots,b_s;w)$ is defined by
$\sum_{q\geq0}w^q/[q!\prod_r(b_r)_q]$. The ratio of successive
terms has magnitude $|w|/[(q+1)\prod_r|q+b_r|]$, which tends to
zero. Its series is entire in $w$, so \eqref{a-spectral-hypergeometric-lane}
is an everywhere convergent identity with the factorial series.
For example $\mathcal A_{3,1}=\{2/3,4/3\}$ and
$E_{3,1}(z)=z\,{}_0F_2(;2/3,4/3;z^3/27)$.

\subsection{The cyclic companion system}

Termwise differentiation on compact subsets gives
$E'_{L,j}=E_{L,j-1}$ for $j\geq1$. For $j=0$, the constant term
vanishes under differentiation and reindexing gives
$E'_{L,0}=E_{L,L-1}$. Let
$\mathbf E_L=(E_{L,0},\ldots,E_{L,L-1})^{\mathsf T}$ and
define the cyclic matrix by $(C_Lv)_0=v_{L-1}$ and
$(C_Lv)_j=v_{j-1}$ for $j\geq1$. Then
\begin{equation}
 \mathbf E_L'=C_L\mathbf E_L,\qquad
 \mathbf E_L(0)=e_0,\qquad
 C_L^L=I,\qquad
 \mathbf E_L(z)=e^{zC_L}e_0.
 \label{a-spectral-companion-system}
\end{equation}
One can verify the last equality directly: $C_L^ne_0=e_{n\bmod L}$,
so its $j$th exponential coordinate is exactly
\eqref{a-spectral-factorial-lane}. Character filtering also gives
\begin{equation}
 E_{L,j}(z)=\frac1L\sum_{\ell=0}^{L-1}
               \omega_L^{-j\ell}e^{\omega_L^\ell z}.
 \label{a-spectral-root-exponential}
\end{equation}
Thus factorial, hypergeometric, companion-state, and root-filter forms
describe the same analytic function. Each form has a different numerical
conditioning: a root sum can contain cancellations even when the
defining series has positive terms.

For a general finite transfer matrix $M_1$, the residue-state identity
$v_{Lq+j}=M_1^j(M_1^L)^qv_0$ follows by adding exponents. It needs
no hypergeometric coefficients. A factorial lane provides a particular
analytic coordinate system; it does not impose a unique microscopic
Hamiltonian on every such transfer problem.

\Needspace{24\baselineskip}
\subsection{Bounded operators and controlled series errors}

\begin{theorem}[Entire lane calculus for a bounded operator]
\label{a-spectral-bounded-operator-lanes}
Let $A_0\in\mathcal B(\HH)$ and $t\in\C$. The series
$E_{L,j}(tA_0)$ converges absolutely in operator norm and satisfies
\begin{align}
 E_{L,j}(tA_0)
 &=\frac1L\sum_{\ell=0}^{L-1}\omega_L^{-j\ell}
                           e^{\omega_L^\ell tA_0},\nonumber\\
 &=\frac{(tA_0)^j}{j!}\,
 {}_0F_{L-1}\!\left(\begin{matrix}-\\\mathcal A_{L,j}\end{matrix};
                      \frac{(tA_0)^L}{L^L}\right),\nonumber\\
 \sum_{j=0}^{L-1}E_{L,j}(tA_0)&=e^{tA_0}.
 \label{a-spectral-operator-lane-identities}
\end{align}
If all terms through degree $Q$ are retained in a lane, the omitted
operator norm is at most
\begin{equation}
 e^a\frac{a^{Q+1}}{(Q+1)!},\qquad a=|t|\norm{A_0}.
 \label{a-spectral-lane-tail-bound}
\end{equation}
\end{theorem}
\begin{proof}
The sum of norms of all exponential terms is at most $e^a$. A lane
is a subseries of this sum, so it converges absolutely. The character
calculation and factorial split can therefore be applied in operator
norm, giving the first two identities. Partitioning the nonnegative
degrees by their residues proves the third. All products are powers
of the same operator, so there is no commuting-operator assumption.
For the tail, bound a lane by the full tail. Since
$(Q+1+k)!\geq(Q+1)!k!$,
\[
 \sum_{n=Q+1}^\infty\frac{a^n}{n!}
 \leq\frac{a^{Q+1}}{(Q+1)!}\sum_{k=0}^\infty\frac{a^k}{k!},
\]
which proves \eqref{a-spectral-lane-tail-bound}.
\end{proof}

If $A_0^L=\alpha I$, the hypergeometric expression reduces every
lane to a scalar function multiplying $A_0^j$. The propagator lies
in the finite span of $I,A_0,\ldots,A_0^{L-1}$. A more general
polynomial identity gives the corresponding finite algebra of powers.
These statements follow from the supplied algebraic relation; they
do not assert one for an arbitrary Hamiltonian.

The coefficient maps $\mathfrak R_{L,j}$ in
\eqref{a-spectral-analytic-lanes} are idempotents on analytic functions.
The evaluated operators $E_{L,j}(tA_0)$ generally are not.
For $L=2$, $E_{2,0}(tA_0)=\cosh(tA_0)$ and
$E_{2,1}(tA_0)=\sinh(tA_0)$, which are neither an orthogonal
projector family nor a pair of independent unitary evolutions.
The sum reconstructs an evolution; an individual lane is an analytic
contribution to it.

\subsection{A squared-factorial lane for Bessel functions}

The same separation gives the entire function
\[
 B(y)=\sum_{n=0}^\infty\frac{(-1)^ny^n}{(n!)^2}
     ={}_0F_1(;1;-y).
\]
For $y\geq0$, this is $J_0(2\sqrt y)$; the entire series itself
does not need a square-root branch. Define
$B_{L,j}(y)=\sum_{q\geq0}(-1)^{Lq+j}y^{Lq+j}/[(Lq+j)!]^2$.
Squaring \eqref{a-spectral-factorial-split} gives
\begin{equation}
 B_{L,j}(y)
 =\frac{(-1)^jy^j}{(j!)^2}\,
 {}_0F_{2L-1}\!\left(
 \begin{matrix}-\\1,\mathcal A_{L,j},\mathcal A_{L,j}\end{matrix};
 \frac{(-1)^Ly^L}{L^{2L}}\right).
 \label{a-spectral-bessel-lane}
\end{equation}
One of the two $q!$ factors is the conventional hypergeometric
factorial and the other supplies the explicit bottom parameter $1$.
All bottom parameters are positive and the ratio test again gives
entire convergence. Character filtering reconstructs $B=\sum_jB_{L,j}$.
With the Euler derivative $\Theta=y\,d/dy$, termwise differentiation
also gives the cyclic system
\begin{equation}
 \Theta^2B_{L,j}(y)=-yB_{L,j-1\bmod L}(y).
 \label{a-spectral-bessel-euler}
\end{equation}
For a term of positive degree $n$, its new coefficient is
$n^2(-1)^n/(n!)^2=-(-1)^{n-1}/[(n-1)!]^2$.
The degree-zero term disappears under $\Theta^2$, which also verifies
the wraparound case $j=0$. The scalar Gowdy Bessel mode therefore
admits these analytic coordinates. Filtering that already-defined
solution does not add a gravitational field equation.



\section{Spectral domains and number-residue packets}
\label{a-spectral-domains-packets}

The bounded calculus has two useful continuations. An unbounded
Hamiltonian admits spectral functions with their exact domains.
An oscillator admits actual residue projectors on its number spectrum.
These are distinct constructions, and only the latter is a family of
orthogonal quantum packet projectors.

\subsection{An unbounded propagator does not make every lane bounded}

Let $H_\infty$ be self-adjoint with spectral measure $E$. For real $t$
define the scalar functions
$f_j(\lambda)=E_{L,j}(-it\lambda)$.
The operator $f_j(H_\infty)$ has domain
\begin{equation}
 \Dom(f_j(H_\infty))
 =\left\{\psi:\int_{\R}|f_j(\lambda)|^2
                       \,d\langle\psi,E(\lambda)\psi\rangle<\infty\right\}.
 \label{a-spectral-lane-domain}
\end{equation}
The scalar identity $\sum_jf_j(\lambda)=e^{-it\lambda}$ gives
\begin{equation}
 \sum_jf_j(H_\infty)\psi=e^{-itH_\infty}\psi,
 \qquad \psi\in\bigcap_j\Dom(f_j(H_\infty)).
 \label{a-spectral-unbounded-reconstruction}
\end{equation}
The equality holds on the stated common domain because it is a finite
sum of spectral functions there. It does not say that each lane acts
on every vector.

For an explicit example, take $L=4$. The scalar formula gives
\[
 f_0(\lambda)=\tfrac12\bigl(\cosh(t\lambda)+\cos(t\lambda)\bigr),
 \quad
 f_2(\lambda)=\tfrac12\bigl(\cos(t\lambda)-\cosh(t\lambda)\bigr).
\]
If $H_\infty$ is multiplication by $\lambda$ on $L^2(\R)$ and
$t\ne0$, these functions define unbounded operators. Their sum
cancels the growing hyperbolic contribution. The full propagator
$e^{-itH_\infty}$ is nevertheless unitary on all of $L^2(\R)$.
For $L=2$, by contrast, the two Schr\"odinger lanes are
$\cos(tH_\infty)$ and $-i\sin(tH_\infty)$, both bounded. A larger
lane count must not inherit this special dyadic boundedness without proof.

A sufficient common-domain condition is
\begin{equation}
 \int_{\R} e^{2|t||\lambda|}
              \,d\langle\psi,E(\lambda)\psi\rangle<\infty,
 \label{a-spectral-exponential-domain}
\end{equation}
since $|f_j(\lambda)|\leq e^{|t||\lambda|}$. Spectral cutoffs
$P_R=E([-R,R])$ supply a dense set satisfying this condition. On
$\Ran P_R$, the restricted Hamiltonian is bounded, the norm-series
identities apply, and
\[
 \sum_jf_j(H_\infty)P_R\psi=e^{-itH_\infty}P_R\psi
      \longrightarrow e^{-itH_\infty}\psi.
\]
On the common domain of the individual lanes, each term also converges
to $f_j(H_\infty)\psi$ by its spectral-domain integrability.
For an arbitrary vector only convergence of the summed cutoff expression
is asserted. The domain question for an interacting Hamiltonian therefore
remains even when a bounded surrogate has excellent numerical residuals.

\subsection{A short factorial tail and its adic unit window}

Let $p$ be prime, let $L=mp$ for a positive integer $m$, and take
$0\leq j<p$. For every $q\geq0$,
\begin{equation}
 g_q=\frac{j!(Lq)!}{(Lq+j)!}
       =\binom{Lq+j}{j}^{-1}\in\Z_p^\times.
 \label{a-spectral-unit-window}
\end{equation}
Indeed
$\binom{Lq+j}{j}=\prod_{r=1}^j(Lq+r)/r$.
Each numerator and denominator is a $p$-adic unit because
$Lq\equiv0\pmod p$ and $1\leq r<p$. For $j=0$ the empty product
is one. The proof also covers $p=2$ in its short window. Outside this
window, a factor can cross a multiple of $p$ and the same proof gives
no unit guarantee. This rational coefficient statement supplies no
claim about the $p$-adic convergence of the complex exponential series.

Analytic refinement of a factorial lane is independent of this unit
condition. Splitting $q=pr+a$, $0\leq a<p$, gives the exact identity
\begin{equation}
 E_{L,j}(z)=\sum_{a=0}^{p-1}E_{pL,j+aL}(z),\qquad 0\leq j<L.
 \label{a-spectral-analytic-refinement}
\end{equation}
Absolute convergence permits this partition of the summation index.
The labels $j+aL$ all lie in $[0,pL)$, but need not lie in the
parent's short unit window. Exact complex analytic refinement and
the local $p$-adic unit assertion therefore have separate hypotheses.

\subsection{Actual orthogonal packets on the oscillator number spectrum}

On the oscillator Fock space, let
$\widehat N|n\rangle=n|n\rangle$ and define the bounded unitary
$V=e^{2\pi i\widehat N/L}$. Its spectral residues are
\begin{equation}
 Q_j^{(L)}=\frac1L\sum_{\ell=0}^{L-1}\omega_L^{-j\ell}V^\ell
          =\sum_{n\equiv j\;({\rm mod}\;L)}|n\rangle\langle n|.
 \label{a-spectral-number-projectors}
\end{equation}
On every number basis vector, the same finite character sum proves the
second equality. These $Q_j^{(L)}$ are bounded orthogonal projectors,
$Q_jQ_k=\delta_{jk}Q_j$, and $\sum_jQ_j=I$. No exponential norm series
for the unbounded number operator is needed; its spectral calculus
defines $V$.

For a coherent state
\[
 |\alpha\rangle=e^{-|\alpha|^2/2}
        \sum_{n=0}^\infty\frac{\alpha^n}{\sqrt{n!}}|n\rangle,
 \qquad \mu=|\alpha|^2\geq0,
\]
the packet probability is the Poisson residue sum
\begin{align}
 p_{L,j}(\mu)
 &=\langle\alpha,Q_j^{(L)}\alpha\rangle
   =e^{-\mu}E_{L,j}(\mu)\nonumber\\
 &=\frac1L\sum_{\ell=0}^{L-1}\omega_L^{-j\ell}
                      e^{\mu(\omega_L^\ell-1)}.
 \label{a-spectral-coherent-packets}
\end{align}
The factorial series proves nonnegativity, including the limiting values
$p_{L,0}(0)=1$ and $p_{L,j}(0)=0$ for $j>0$. Partition of the
Poisson series proves $\sum_jp_{L,j}=1$. Their equivalent hypergeometric
form is $e^{-\mu}$ times \eqref{a-spectral-hypergeometric-lane} at $z=\mu$.

\begin{proposition}[Approach to uniform residue probability]
\label{a-spectral-coherent-bound}
For fixed $L\geq2$ and $\mu\geq0$,
\begin{equation}
 \left|p_{L,j}(\mu)-\frac1L\right|
 \leq\frac1L\sum_{\ell=1}^{L-1}
       e^{-\mu(1-\cos(2\pi\ell/L))}
 \leq\frac{L-1}{L}
       e^{-\mu(1-\cos(2\pi/L))}.
 \label{a-spectral-packet-mixing-bound}
\end{equation}
\end{proposition}
\begin{proof}
In \eqref{a-spectral-coherent-packets}, the $\ell=0$ term is $1/L$.
The modulus of every other term is
$e^{\mu(\operatorname{Re}\omega_L^\ell-1)}/L$.
The triangle inequality gives the first bound. Among the nonzero
residues, the largest cosine is $\cos(2\pi/L)$, giving the second.
\end{proof}

The exponent is positive for each fixed $L$, but decreases as
$2\pi^2/L^2$ when $L$ grows. The conclusion is therefore uniform in
the finite packet label $j$ at fixed $L$; it is not a depth-uniform
estimate at arbitrarily large residue modulus.

Finally the annihilation operator routes adjacent residues:
\begin{equation}
 Q_r^{(L)}aQ_s^{(L)}=0\quad\text{unless }r\equiv s-1\pmod L,
 \qquad
 Q_r^{(L)}a^\dagger Q_s^{(L)}=0\quad\text{unless }r\equiv s+1\pmod L.
 \label{a-spectral-ladder-routing}
\end{equation}
On a finite number-state combination this follows from
$a|n\rangle=\sqrt n\,|n-1\rangle$ and
$a^\dagger|n\rangle=\sqrt{n+1}\,|n+1\rangle$.
Each residue projector preserves $\Dom(a)$ and $\Dom(a^\dagger)$,
because it only removes terms from their weighted number norm.
Truncation of a vector in either domain converges in the corresponding
graph norm, so the identities extend to those domains.
For the isolated oscillator Hamiltonian, which is a function of
$\widehat N$, the packet projectors commute with the evolution and
the weights stay constant. A driving term containing ladder operators
can transfer weight according to the routing rule. Selecting such a
drive is further model data.

These packet probabilities are directly interpretable within the
ordinary oscillator model. A physical interpretation of adic digits,
hidden tree levels, or a preferred hypergeometric residue split requires
an independently specified coupling and observable. The exact analytic
and spectral coordinates derived here remain useful without making
that identification.

